\chapter{The Two-Dimensional Classical Ising Model}
\label{chap:two-dimensional-classical-ising}

The two-dimensional classical Ising model is the first genuinely nontrivial
exactly solved lattice model in these notes.  It is also the central bridge
between two themes developed earlier: transfer matrices and free fermions.  In
one dimension, the classical Ising transfer matrix is only a $2\times2$ matrix;
in two dimensions, a row-to-row transfer matrix acts on a $2^L$-dimensional row
Hilbert space.  The remarkable fact, discovered by Onsager and reformulated by
Kaufman, is that this exponentially large transfer matrix is still exactly
solvable because it is a Gaussian operator in Majorana variables.

This chapter has three goals.  First, we define the anisotropic zero-field model
and construct its row-to-row transfer matrix in the same convention used in the
quantum--classical correspondence chapter.  Second, we explain why the transfer
matrix belongs to a free-fermion algebra.  Third, we extract the Onsager critical
condition
\begin{equation}
\label{eq:2d-classical-ising-onsager-critical-preview}
    \sinh(2K_x)\sinh(2K_y)=1
\end{equation}
from the closing of the transfer-matrix gap.

\begin{remark}[Zero field is essential]
The exact free-fermion solution discussed here applies to the square-lattice
Ising model at zero magnetic field.  Adding a field term
$H\sum_{j,m}s_{j,m}$ destroys the Gaussian structure of the transfer matrix.
The two-dimensional Ising model in a nonzero magnetic field is not known to have
a comparable closed-form exact solution.
\end{remark}

\section{Classical partition function}
\label{sec:2d-classical-ising-partition-function}

Consider an $L\times M$ square lattice with periodic boundary conditions in both
directions.  A classical spin variable
\begin{equation}
    s_{j,m}=\pm1,
    \qquad
    j=1,\ldots,L,
    \qquad
    m=1,\ldots,M,
\end{equation}
lives at each lattice site.  We identify
\begin{equation}
    s_{L+1,m}=s_{1,m},
    \qquad
    s_{j,M+1}=s_{j,1}.
\end{equation}
The anisotropic zero-field Ising partition function is
\begin{equation}
\label{eq:2d-classical-ising-partition-function}
    Z_{L,M}(K_x,K_y)
    =
    \sum_{\{s_{j,m}=\pm1\}}
    \exp\left[
        K_x\sum_{j,m}s_{j,m}s_{j+1,m}
        +
        K_y\sum_{j,m}s_{j,m}s_{j,m+1}
    \right].
\end{equation}
Here $K_x$ and $K_y$ are dimensionless couplings.  If the physical Hamiltonian is
\begin{equation}
    E[s]
    =
    -J_x\sum_{j,m}s_{j,m}s_{j+1,m}
    -J_y\sum_{j,m}s_{j,m}s_{j,m+1},
\end{equation}
then
\begin{equation}
    K_x=\beta J_x,
    \qquad
    K_y=\beta J_y .
\end{equation}
For most of this chapter we assume $K_x,K_y>0$, corresponding to the
ferromagnetic model.  Antiferromagnetic signs can often be moved by sublattice
spin flips on a bipartite lattice, but the ferromagnetic case is the canonical
one for discussing the thermal order--disorder transition.

Let
\begin{equation}
    \bm{s}_m=(s_{1,m},\ldots,s_{L,m})
\end{equation}
denote the $m$th row configuration.  The Boltzmann weight between two adjacent
rows can be distributed symmetrically as
\begin{equation}
\label{eq:2d-classical-ising-symmetric-row-weight}
    T_{\bm{s},\bm{s}'}
    =
    \exp\left[
        \frac{K_x}{2}\sum_{j=1}^{L}s_js_{j+1}
        +
        K_y\sum_{j=1}^{L}s_js_j'
        +
        \frac{K_x}{2}\sum_{j=1}^{L}s_j's_{j+1}'
    \right].
\end{equation}
Then
\begin{equation}
\label{eq:2d-classical-ising-z-tr-tm}
    Z_{L,M}(K_x,K_y)=\Tr T^M .
\end{equation}
The symmetrization in \cref{eq:2d-classical-ising-symmetric-row-weight} is not
physically necessary, but it makes $T$ real symmetric and therefore convenient
for spectral analysis.

\section{Operator form of the row-to-row transfer matrix}
\label{sec:2d-classical-ising-row-transfer-operator}

We now rewrite $T$ as an operator acting on a row Hilbert space
\begin{equation}
    \calH_{\rm row}=(\CC^2)^{\otimes L}.
\end{equation}
To remain consistent with the transverse-field Ising convention used earlier in
these notes, we label a classical row configuration by eigenvalues of
$\sigma_j^x$:
\begin{equation}
\label{eq:2d-classical-ising-x-basis}
    \sigma_j^x\ket{s_1,\ldots,s_L}_x
    =
    s_j\ket{s_1,\ldots,s_L}_x .
\end{equation}
In this basis, the intra-row interaction is diagonal:
\begin{equation}
\label{eq:2d-classical-ising-horizontal-operator}
    \exp\left(
        \frac{K_x}{2}\sum_{j=1}^{L}\sigma_j^x\sigma_{j+1}^x
    \right).
\end{equation}
The vertical one-site Boltzmann matrix is
\begin{equation}
\label{eq:2d-classical-ising-vertical-weight}
    W_{ss'}=\ee^{K_yss'}
    =
    \begin{pmatrix}
        \ee^{K_y} & \ee^{-K_y} \\
        \ee^{-K_y} & \ee^{K_y}
    \end{pmatrix}_{s,s'} .
\end{equation}
Since $\sigma^z$ flips the $\sigma^x$ eigenvalue, this matrix can be written as
\begin{equation}
\label{eq:2d-classical-ising-vertical-exp}
    W=A\ee^{\kappa\sigma^z},
    \qquad
    \tanh\kappa=\ee^{-2K_y},
    \qquad
    A=\sqrt{2\sinh(2K_y)} .
\end{equation}
Indeed,
\begin{equation}
    A\cosh\kappa=\ee^{K_y},
    \qquad
    A\sinh\kappa=\ee^{-K_y}.
\end{equation}
The parameter $\kappa$ is the dual coupling associated with $K_y$:
\begin{equation}
\label{eq:2d-classical-ising-kappa-dual}
    \ee^{-2K_y}=\tanh\kappa
    \quad\Longleftrightarrow\quad
    \sinh(2K_y)\sinh(2\kappa)=1 .
\end{equation}
Thus the symmetric row-to-row transfer matrix is
\begin{equation}
\label{eq:2d-classical-ising-symmetric-transfer-operator}
    T
    =
    A^L
    \exp\left(
        \frac{K_x}{2}\sum_{j=1}^{L}\sigma_j^x\sigma_{j+1}^x
    \right)
    \exp\left(
        \kappa\sum_{j=1}^{L}\sigma_j^z
    \right)
    \exp\left(
        \frac{K_x}{2}\sum_{j=1}^{L}\sigma_j^x\sigma_{j+1}^x
    \right).
\end{equation}
This is the precise point at which the two-dimensional classical model starts to
look like a one-dimensional quantum spin chain: $T$ is an exponential product of
operators acting on a one-dimensional row Hilbert space.

\begin{remark}[Relation to the notation $K_\tau$]
In the quantum--classical correspondence chapter, the transfer direction was
written as an imaginary-time direction and the vertical coupling was denoted
$K_\tau$.  In the present chapter we write the two classical directions as
$x$ and $y$.  Thus $K_y$ here plays the role of $K_\tau$ there.
\end{remark}

\section{Kaufman's free-fermion structure}
\label{sec:2d-classical-ising-kaufman-free-fermion}

At first sight, \cref{eq:2d-classical-ising-symmetric-transfer-operator} is not
easier than the original classical problem, because it acts on a
$2^L$-dimensional vector space.  The key observation is that both exponential
factors are Gaussian after the Jordan--Wigner transformation.

Using the Jordan--Wigner convention fixed earlier in these notes, define
Majorana operators
\begin{equation}
\label{eq:2d-classical-ising-majorana-definitions}
    a_{2j-1}
    =
    \left(\prod_{\ell<j}\sigma_\ell^z\right)\sigma_j^x,
    \qquad
    a_{2j}
    =
    \left(\prod_{\ell<j}\sigma_\ell^z\right)\sigma_j^y .
\end{equation}
They satisfy
\begin{equation}
\label{eq:2d-classical-ising-majorana-algebra}
    a_m^\dagger=a_m,
    \qquad
    \{a_m,a_n\}=2\delta_{mn} .
\end{equation}
In the bulk one has the local identities
\begin{equation}
\label{eq:2d-classical-ising-majorana-spin-identities}
    \sigma_j^z=-\ii a_{2j-1}a_{2j},
    \qquad
    \sigma_j^x\sigma_{j+1}^x=-\ii a_{2j}a_{2j+1} .
\end{equation}
Therefore the two building blocks of $T$ become
\begin{align}
\label{eq:2d-classical-ising-vertical-majorana-gaussian}
    \exp\left(\kappa\sum_j\sigma_j^z\right)
    &=
    \exp\left(-\ii\kappa\sum_j a_{2j-1}a_{2j}\right),
    \\
\label{eq:2d-classical-ising-horizontal-majorana-gaussian}
    \exp\left(\frac{K_x}{2}\sum_j\sigma_j^x\sigma_{j+1}^x\right)
    &=
    \exp\left(-\frac{\ii K_x}{2}\sum_j a_{2j}a_{2j+1}\right),
\end{align}
up to the usual boundary-sector term for a periodic spin chain.

Equations \cref{eq:2d-classical-ising-vertical-majorana-gaussian,eq:2d-classical-ising-horizontal-majorana-gaussian}
show that $T$ is a product of exponentials of Majorana bilinears.  Such operators
are fermionic Gaussian operators.  Their adjoint action maps Majoranas linearly
to Majoranas:
\begin{equation}
\label{eq:2d-classical-ising-gaussian-adjoint-action}
    T^{-1}a_mT
    =
    \sum_{n=1}^{2L}R_{mn}a_n,
\end{equation}
where $R$ is a $2L\times2L$ matrix preserving the Majorana anticommutation
relations.  The exponentially large transfer-matrix problem is therefore reduced
to the diagonalization of a one-particle Majorana transformation.

\begin{remark}[Boundary sectors]
For periodic spin boundary conditions, the boundary term in
$\sigma_L^x\sigma_1^x$ contains the fermion-parity operator
$P_F=\prod_{j=1}^{L}\sigma_j^z$.  Consequently, the exact finite-$L$ partition
function is obtained by combining parity-projected fermion traces with periodic
and antiperiodic momenta, in the same spirit as the TFIM parity-sector
bookkeeping.  In the thermodynamic limit these sector distinctions affect only
subleading finite-size terms, not the bulk free energy or the critical curve.
\end{remark}

Because the model is translation invariant, the one-particle Majorana problem
block-diagonalizes in momentum.  For each momentum $k$, the relevant
$2\times2$ block has eigenvalues $\ee^{\pm\varepsilon(k)}$, where
\begin{equation}
\label{eq:2d-classical-ising-onsager-dispersion}
    \cosh\varepsilon(k)
    =
    \cosh(2\kappa)\cosh(2K_x)
    -
    \sinh(2\kappa)\sinh(2K_x)\cos k .
\end{equation}
The quantity $\varepsilon(k)$ is the transfer-matrix analogue of a quasiparticle
energy.  More precisely, it is the single-particle gap in the logarithm of the
transfer matrix.  Its minimum occurs at $k=0$ for ferromagnetic couplings, and
\begin{equation}
\label{eq:2d-classical-ising-transfer-gap}
    \varepsilon(0)
    =
    2\abs{\kappa-K_x} .
\end{equation}
Thus the transfer-matrix gap closes when
\begin{equation}
\label{eq:2d-classical-ising-kappa-equals-kx}
    \kappa=K_x .
\end{equation}
Using \cref{eq:2d-classical-ising-kappa-dual}, this is equivalent to
\begin{equation}
\label{eq:2d-classical-ising-onsager-critical}
    \boxed{\sinh(2K_x)\sinh(2K_y)=1.}
\end{equation}
This is the Onsager critical line for the anisotropic two-dimensional Ising
model.

\section{Largest eigenvalue and Onsager free energy}
\label{sec:2d-classical-ising-onsager-free-energy}

Let
\begin{equation}
\label{eq:2d-classical-ising-phi-definition}
    \Phi(K_x,K_y)
    =
    \lim_{L,M\to\infty}\frac{1}{LM}\log Z_{L,M}(K_x,K_y)
\end{equation}
be the dimensionless free-energy density with the sign convention appropriate to
$\log Z$.  The physical free energy per site is
\begin{equation}
    f=-\beta^{-1}\Phi
\end{equation}
when $K_\mu=\beta J_\mu$.  Since $Z=\Tr T^M$, $\Phi$ is the logarithm of the
largest transfer-matrix eigenvalue per row site.

From the Gaussian diagonalization one obtains
\begin{equation}
\label{eq:2d-classical-ising-onsager-free-energy-kappa}
    \Phi(K_x,K_y)
    =
    \frac12\log\!\bigl(2\sinh(2K_y)\bigr)
    +
    \frac12\int_{-\pi}^{\pi}\frac{\dd k}{2\pi}\,\varepsilon(k),
\end{equation}
where $\varepsilon(k)$ is defined by
\cref{eq:2d-classical-ising-onsager-dispersion}.  Equivalently,
\begin{align}
\label{eq:2d-classical-ising-onsager-free-energy-standard}
    \Phi(K_x,K_y)
    &=
    \frac12\log\!\bigl(2\sinh(2K_y)\bigr)
    +
    \frac{1}{2\pi}
    \int_0^\pi \dd k\,\operatorname{arcosh} X(k),
    \\
    X(k)
    &=
    \cosh(2\kappa)\cosh(2K_x)
    -
    \sinh(2\kappa)\sinh(2K_x)\cos k .
\end{align}
The appearance of $\operatorname{arcosh}$ is the transfer-matrix version of
``taking a quasiparticle energy from a one-particle evolution eigenvalue.''

A useful check is the decoupled-row limit $K_x=0$.  Then
\begin{equation}
    \varepsilon(k)=2\kappa,
\end{equation}
and \cref{eq:2d-classical-ising-onsager-free-energy-kappa} gives
\begin{equation}
    \Phi(0,K_y)
    =
    \frac12\log\!\bigl(2\sinh(2K_y)\bigr)+\kappa
    =
    \log(2\cosh K_y),
\end{equation}
which is exactly the free-energy density of independent vertical Ising chains
viewed through the chosen transfer direction.

For the isotropic model $K_x=K_y=K$, the critical point follows from
\begin{equation}
    \sinh^2(2K_c)=1,
\end{equation}
and therefore
\begin{equation}
\label{eq:2d-classical-ising-isotropic-kc}
    \boxed{K_c=\frac12\log(1+\sqrt2).}
\end{equation}
If $K=\beta J$, this gives
\begin{equation}
\label{eq:2d-classical-ising-isotropic-tc}
    k_B T_c
    =
    \frac{2J}{\log(1+\sqrt2)} .
\end{equation}

\section{Duality interpretation of the critical curve}
\label{sec:2d-classical-ising-duality-interpretation}

The same critical condition follows from Kramers--Wannier duality.  The
high-temperature expansion writes each bond factor as
\begin{equation}
\label{eq:2d-classical-ising-high-temp-bond}
    \ee^{Ks_rs_{r'}}
    =
    \cosh K\left(1+s_rs_{r'}\tanh K\right).
\end{equation}
After summing over spins, only closed even-degree bond configurations survive.
The low-temperature expansion of a dual Ising model is instead a sum over domain
walls.  Matching high-temperature loops of the original lattice with
domain-wall loops of the dual lattice gives
\begin{equation}
\label{eq:2d-classical-ising-dual-couplings}
    \ee^{-2K_x^*}=\tanh K_y,
    \qquad
    \ee^{-2K_y^*}=\tanh K_x .
\end{equation}
Equivalently,
\begin{equation}
\label{eq:2d-classical-ising-dual-couplings-sinh}
    \sinh(2K_x^*)\sinh(2K_y)=1,
    \qquad
    \sinh(2K_y^*)\sinh(2K_x)=1 .
\end{equation}
A self-dual point satisfies
\begin{equation}
    K_x^*=K_x,
    \qquad
    K_y^*=K_y,
\end{equation}
and hence obeys
\begin{equation}
    \sinh(2K_x)\sinh(2K_y)=1 .
\end{equation}
Duality alone identifies the possible transition curve if there is a unique
ferromagnetic critical line.  The transfer-matrix solution gives the sharper
statement: the one-particle transfer gap actually closes on this curve.

\section{Order parameter and critical behavior}
\label{sec:2d-classical-ising-order-parameter-critical-behavior}

The finite-volume partition function at zero field has a global $\ZZ_2$ symmetry
$s_{j,m}\mapsto -s_{j,m}$, so the finite-volume magnetization vanishes exactly.
The spontaneous magnetization is defined by taking the thermodynamic limit before
removing an infinitesimal symmetry-breaking field:
\begin{equation}
\label{eq:2d-classical-ising-spontaneous-magnetization-definition}
    M
    =
    \lim_{H\to0^+}\lim_{L,M\to\infty}
    \frac{1}{LM}\sum_{j,m}\langle s_{j,m}\rangle_H .
\end{equation}
Yang's exact result for the anisotropic square-lattice model is
\begin{equation}
\label{eq:2d-classical-ising-yang-magnetization}
    M(K_x,K_y)
    =
    \begin{cases}
    \displaystyle
    \left[
        1-\frac{1}{\sinh^2(2K_x)\sinh^2(2K_y)}
    \right]^{1/8},
    &
    \sinh(2K_x)\sinh(2K_y)>1,\\[12pt]
    0,
    &
    \sinh(2K_x)\sinh(2K_y)\leq1 .
    \end{cases}
\end{equation}
For the isotropic model this reduces to
\begin{equation}
\label{eq:2d-classical-ising-isotropic-magnetization}
    M(K)
    =
    \left[1-\sinh^{-4}(2K)\right]^{1/8},
    \qquad
    K>K_c .
\end{equation}
Thus the magnetization exponent is
\begin{equation}
\label{eq:2d-classical-ising-beta-exponent}
    M\sim (T_c-T)^{1/8},
    \qquad
    \beta_{\rm Ising}=\frac18 .
\end{equation}

The transfer-matrix dispersion also reveals the divergence of the correlation
length.  Along the transfer direction, the inverse correlation length is given by
the relevant transfer gap.  Near the ferromagnetic critical line,
\begin{equation}
\label{eq:2d-classical-ising-correlation-length-gap}
    \xi_y^{-1}\sim\varepsilon(0)=2\abs{\kappa-K_x} .
\end{equation}
Because $\kappa-K_x$ is linear in the distance from the critical curve, the
correlation-length exponent is
\begin{equation}
\label{eq:2d-classical-ising-nu-exponent}
    \xi\sim\abs{T-T_c}^{-1},
    \qquad
    \nu=1 .
\end{equation}
The full critical behavior of the two-dimensional Ising universality class is
summarized by
\begin{equation}
\label{eq:2d-classical-ising-critical-exponents}
\begin{gathered}
    \alpha=0\quad(\text{logarithmic specific-heat divergence}),
    \qquad
    \beta=\frac18,
    \qquad
    \gamma=\frac74,
    \\
    \nu=1,
    \qquad
    \eta=\frac14,
    \qquad
    \delta=15 .
\end{gathered}
\end{equation}
In these notes, the most important point is not the list of exponents itself,
but their origin: the classical thermal transition is controlled by a massless
free-fermion transfer-matrix theory at criticality.

\section{Connection to the one-dimensional TFIM}
\label{sec:2d-classical-ising-connection-to-tfim}

The transfer matrix in
\cref{eq:2d-classical-ising-symmetric-transfer-operator} becomes the imaginary-
time evolution operator of the one-dimensional TFIM in a strongly anisotropic
limit.  Set
\begin{equation}
\label{eq:2d-classical-ising-tfim-scaling}
    K_x=\delta\tau J,
    \qquad
    \kappa=\delta\tau h,
    \qquad
    \tanh(\delta\tau h)=\ee^{-2K_y} .
\end{equation}
Then, using the symmetric Trotter formula,
\begin{equation}
\label{eq:2d-classical-ising-tfim-transfer-limit}
    T
    =
    A^L\exp\left[-\delta\tau H_{\rm TFIM}+O(\delta\tau^3)\right],
\end{equation}
with
\begin{equation}
\label{eq:2d-classical-ising-tfim-hamiltonian}
    H_{\rm TFIM}
    =
    -J\sum_{j=1}^{L}\sigma_j^x\sigma_{j+1}^x
    -h\sum_{j=1}^{L}\sigma_j^z .
\end{equation}
The Onsager condition becomes
\begin{equation}
    \sinh(2\delta\tau J)\sinh(2K_y)=1.
\end{equation}
Using
\begin{equation}
    \ee^{-2K_y}=\tanh(\delta\tau h)
    =
    \delta\tau h+O(\delta\tau^3),
\end{equation}
one finds
\begin{equation}
\label{eq:2d-classical-ising-tfim-critical-limit}
    h=J .
\end{equation}
Thus the thermal critical point of the anisotropic two-dimensional classical
Ising model maps to the zero-temperature quantum critical point of the
one-dimensional transverse-field Ising model.

\begin{center}
\begin{tabular}{@{}L{0.31\linewidth}L{0.31\linewidth}L{0.28\linewidth}@{}}
\toprule
Classical 2D Ising object & Transfer-matrix object & 1D quantum interpretation \\
\midrule
Row configuration $\bm{s}_m$ & Basis vector in $\calH_{\rm row}$ & Spin-chain configuration \\
Row-to-row transfer matrix $T$ & Positive operator on $\calH_{\rm row}$ & Imaginary-time step $\ee^{-\delta\tau H}$ \\
Largest eigenvalue of $T$ & Ground-state sector of $-\log T$ & Free-energy density / ground-state energy density \\
Transfer gap $\varepsilon(0)$ & Mass gap of $-\log T$ & Quantum excitation gap \\
Onsager curve & Gapless transfer matrix & TFIM critical point $h=J$ \\
\bottomrule
\end{tabular}
\end{center}

\section{Summary}
\label{sec:2d-classical-ising-summary}

The two-dimensional zero-field Ising model is exactly solvable because its
row-to-row transfer matrix is a fermionic Gaussian operator.  The essential chain
of ideas is
\begin{equation}
\label{eq:2d-classical-ising-solution-chain}
\begin{gathered}
    \text{2D classical partition function}
    \longrightarrow
    \text{row transfer matrix}
    \longrightarrow
    \text{Majorana Gaussian operator}
    \\
    \longrightarrow
    \text{one-particle dispersion }\varepsilon(k).
\end{gathered}
\end{equation}
The Onsager critical point is the point where the transfer-matrix quasiparticle
gap closes:
\begin{equation}
    \varepsilon(0)=0
    \quad\Longleftrightarrow\quad
    \kappa=K_x
    \quad\Longleftrightarrow\quad
    \sinh(2K_x)\sinh(2K_y)=1 .
\end{equation}
This is why the model belongs naturally in a lecture note on solvable quantum
many-body systems: its exact solution is not merely a classical combinatorial
miracle, but a free-fermion solution in transfer-matrix language.
